%%% Einfaches Template für einen Abschlussbericht zum Berufspraktikum

%%% Magic Comments zum Setzen der korrekten Parameter in kompatiblen IDEs
% !TeX encoding = utf8
% !TeX program = pdflatex 
% !TeX spellcheck = de_DE
% !BIB program = biber

\RequirePackage[utf8]{inputenc} % bei Verw. von lualatex oder xelatex entfernen!
\RequirePackage{hgbpdfa}				% Remove if NO PDF/A output is desired (fonts are not affected)

\documentclass[type=internship,theme=default,language=german,titlelanguage=german,smartquotes=true]{hgbthesis}
% Zulässige Optionen in [..]:
%    Typ der Arbeit (type=): master (default), bachelor, diploma, phd, internship
%    Theme der Titelseite (theme=): default (default), fhooe26, classic, fhooe24
%    Als Exposé verwenden: proposal=true, proposal=false (default)
%    Hauptsprache im Dokument (language=): german, english (default)
%    Sprache der Titelseite (titlelanguage=): german, english (default: Hauptsprache)
%    Typografische Anführungszeichen: smartquotes=true, smartquotes=false (default)
%    APA Zitierstil: apa=true, apa=false (default)
%    Layout: twoside=false (einseitig, default), twoside=true (zweiseitig)
%    Update-Check: updatecheck=true (default), updatecheck=false (zum Deaktivieren)
%%%-----------------------------------------------------------------------------

\graphicspath{{images/}}     % Verzeichnis mit Bildern und Grafiken
\bibliography{hgbreferences} % Biblatex-Literaturdatei (references.bib)

%%%-----------------------------------------------------------------------------
\begin{document}
%%%-----------------------------------------------------------------------------

%%%-----------------------------------------------------------------------------
% Angaben für die Titelei (Titelseite, Eidesstattliche Erklärung etc.)
%%%-----------------------------------------------------------------------------

\title{Endbericht zum Berufspraktikum bei Mogulovich International}
%\subtitle{Untertitel nur wenn unbedingt notwendig}
\author{Alex A.\ Schlaumeier}

%\programtype{Bachelorstudiengang} % optional, Standard: abgeleitet aus type=
\programname{Medientechnik und -design}
%\institution{Fachhochschule Oberösterreich} % optional, Standard: Name der Hochschule laut Theme

\placeofstudy{Hagenberg}
\dateofsubmission{2027}{07}{01} % {JJJJ}{MM}{TT}

% Liste der Kontaktpersonen, bis zu 4 sind möglich, Titel in [] ist optional
\advisor{Pjotr I.~Czar, M.A.}
%\advisor[Team Lead]{Bettina B. Ackend}

\companyName{%
   Mogulovich International Media GmbH\\
   Online Division\\
   Hubertusgasse 3a, 1020 Wien
}
\companyUrl{www.mogul.at}

\license{cc}      % Unter Creative Commons Lizenz veröffentlichen (empfohlen)
%\license{strict} % Restriktieve Lizenz, "Alle Rechte vorbehalten"

%%%-----------------------------------------------------------------------------
\frontmatter                                       % Titelei (röm. Seitenzahlen)
%%%-----------------------------------------------------------------------------

\maketitle

% Dokumentation der Nutzung von KI-Tools: zutreffende Textbausteine verwenden, Unzutreffendes
% bitte löschen. Weitere Arten der Nutzung, die auf dieser Liste nicht erscheinen, sind unbedingt
% vorab mit Ihrem/Ihrer Betreuer/in zu klären. Ausführliche Dokumentation (aiusagedetails): siehe Manual.
\begin{aiusage}
	%\aiusednone                          % keinerlei Hilfsmittel der generativen KI benutzt
	%\aiusedfor{ideas}                    % um erste Ideen für Forschungsfragen und Theorien zu sammeln
	\aiusedfor{literature-search}         % um Literatur zu finden
	%\aiusedfor{literature-understanding} % um Literatur besser nachvollziehen zu können
	\aiusedfor{proofreading}              % um Korrektur zu lesen und den sprachlichen Stil zu verbessern
	%\aiusedforother{um ... zu ...}       % weitere Anwendungen
\end{aiusage}

\tableofcontents

%%%-----------------------------------------------------------------------------

\chapter{Kurzfassung}

Umfang der Kurzfassung: ca.\ 200 Worte.

Zum allgemeinen Inhalt des Berichts: Dieser Bericht beschreibt den Ablauf des
Praktikums, die Aufgaben und durchgeführten Projekte und Erfahrungen. Die
eigenen Aktivitäten (Projekte) stehen dabei natürlich im Mittelpunkt und bilden
den Hauptteil des Berichts. Wenn viele Kleinprojekte bearbeitet wurden, sollten
einige davon exemplarisch genauer beschrieben werden. Neben der eigentlichen
Arbeit sollten aber auch folgende weitere Aspekte berücksichtigt werden:
%
\begin{itemize}
	\item Abläufe (Workflows) innerhalb des Unternehmens bzw.\ in Projekten
	(grafische Darstellungen können dabei nützlich sein),
	\item Arbeits- und Führungsstil, Kommunikation innerhalb des Unternehmens,
	\item Kommunikation nach außen (Kund*innen, Partner*innen),
	\item Zeitsituation, Terminprobleme,
	\item Einbettung in das Team, soziale Erfahrungen,
	\item Einsatz von speziellen Techniken, Methoden und Werkzeugen,
	\item wichtige Herausforderungen oder Schwierigkeiten,
	\item Anforderungen in Bezug auf die Ausbildung im Studium (gut
	einsetzbare Kenntnisse, Defizite).
\end{itemize}
%
Die nachfolgenden Kapitelüberschriften sollen nur zur Orientierung für die
Struktur des Berichts dienen, über die konkrete Einteilung und den Wortlaut
kann man natürlich selbst entscheiden.

%%%-----------------------------------------------------------------------------
\mainmatter                             % Hauptteil (ab hier arab. Seitenzahlen)
%%%-----------------------------------------------------------------------------

\chapter{Das Unternehmen}

Umfang: 1--2 Seiten

%%%-----------------------------------------------------------------------------

\chapter{Projekte und Tätigkeiten während des Praktikums}

Umfang: 2--3 Seiten (Projektziel(e), Projektumfeld)

%%%-----------------------------------------------------------------------------
     
\chapter{Projektbeispiele}

Umfang: 5--6 Seiten (Umsetzung, grober Terminplan, Ergebnisse,
Qualitätssicherungsmaßnahmen)

%%%-----------------------------------------------------------------------------

\chapter{Erfahrungen und Zusammenfassung}

Umfang: 1--2 Seiten

%%%-----------------------------------------------------------------------------
%\backmatter        % Schlussteil (nur bei Verwendung des Quellenverzeichnisses)
%%%-----------------------------------------------------------------------------

%\MakeBibliography % Quellenverzeichnis (sofern notwendig, sonst weglassen)

\end{document}
